\documentclass{article}
\usepackage{graphicx}
\usepackage{amsmath}
\usepackage{amsfonts}
\usepackage{bookmark}
\usepackage{booktabs}
\usepackage{exercises}
\graphicspath{ {./images/10.4} }
\usepackage{indentfirst}
\usepackage{amsthm}
\usepackage{yhmath}
\usepackage{geometry}
\usepackage{placeins}
\usepackage{csquotes}
\usepackage{pgfplots}
\newcommand*{\oldneg}{\mathord{\sim}}
\newcommand\showdiv[1]{\overline{\smash{)}#1}}
\newcommand\Mydiv[1]{%
\smash{\raise.3ex\hbox{$\big)$}}\mkern-8mu
        \overline{\enspace\strut#1}}
\setlength{\tabcolsep}{5pt}
\theoremstyle{definition}
\newtheorem{definition}{Definition}

\setlength{\parindent}{.5in}
\setlength{\parskip}{12pt}
\renewcommand{\arraystretch}{2}
\geometry{margin=1in}
\newcommand{\?}{\stackrel{?}{=}}
\pgfplotsset{width=7cm,compat=1.18}\usepgfplotslibrary{statistics}

\title{10.4 Answers}
\begin{document}
\maketitle

\section*{3}
\[ 2+3\cos 4\pi-5\? 0 \]
\[ -3+3\cos 4\pi\? 0 \]
\[ -3+3=0 \]

\section*{4}
\[ \pi\sec\pi+\pi\? 0 \]
\[ \pi(-1)+\pi =0 \]

\section*{5}
\[ 12\sin^2\frac{\pi}{6}-3\? 0 \]
\[ 12(\frac{1}{2})^2-3\? 0 \]
\[ 3-3=0 \]

\section*{6}
\[ 5\tan^3\frac{\pi}{4}-5\? 0 \]
\[ 5(1)^3-5\? 0 \]
\[ 5-5=0 \]

\section*{7}
\[ 2\cos^4\frac{\pi}{2}-\cos^2\frac{\pi}{2}\? 0 \]
\[ 2(0)^4-(0)^2\? 0 \]
\[ 0=0 \]

\section*{9}
\[ 2\sin x-1=0 \]
\[ \sin x=\frac{1}{2} \]
In the interval $0\le x\le 2\pi$, one solution is $\frac{\pi}{6}$, while the other is $\pi-\frac{\pi}{6}=\frac{5\pi}{6}$. Thus the general solution of the equation 
is \[ x=\frac{\pi}{6}+2\pi n\vee x=\frac{5\pi}{6}+2\pi n \] where $n$ is an integer.

\section*{10}
\[ \sqrt{3}\csc x+2=0 \]
\[ \csc x=-\frac{2}{\sqrt{3}} \]
In the interval $0\le x\le 2\pi$, one solution is $2\pi-\frac{\pi}{3}=\frac{5\pi}{3}$ 
while the other is $\frac{\pi}{3}+\pi=\frac{4\pi}{3}$ Thus the general solution of the equation is
\[ x=\frac{5\pi}{3}+2\pi n\vee x=\frac{4\pi}{3}+2\pi n \] where $n$ is an integer.

\section*{11}
\[ 3\tan x-\sqrt{3}=0 \]
\[ \tan x=\frac{\sqrt{3}}{3} \]
\[ \tan x=\frac{1}{\sqrt{3}} \]
In the interval $0\le x\le 2\pi$, one solution is $\frac{\pi}{6}$ while the other is $\frac{\pi}{6}+\pi=\frac{7\pi}{6}$.
Thus the general solution of the equation is
\[ x=\frac{\pi}{6}+\pi n\vee x=\frac{7\pi}{6}+\pi n \] where $n$ is an integer.

\section*{12}
\[ \sin x+\sqrt{2}=-\sin x \]
\[ 2\sin x=-\sqrt{2} \]
\[ \sin x=-\frac{\sqrt{2}}{2} \]
\[ \sin x=-\frac{1}{\sqrt{2}} \]
In the interval $0\le x\le 2\pi$, one solution is $2\pi-\frac{\pi}{4}=\frac{7\pi}{4}$ while the other is $\frac{\pi}{4}+\pi=\frac{5\pi}{4}$.
Thus the general solution to the equation is
\[ x=\frac{7\pi}{4}+2\pi n\vee x=\frac{5\pi}{4}+2\pi n \] where $n$ is an integer.

\section*{13}
\[ 4\cos^2 x-3=0 \]
\[ \cos^2 x=\frac{3}{4} \]
\[ \cos x=\pm\frac{\sqrt{3}}{2} \]
In the interval $0\le x\le 2\pi$, the solutions are $\frac{\pi}{6},\frac{11\pi}{6},\pi-\frac{\pi}{6}=\frac{5\pi}{6},$ and 
$\frac{\pi}{6}+\pi=\frac{7\pi}{6}$. Thus the general solution to the equation is
\[ x=\frac{\pi}{6}+2\pi n\vee x=\frac{11\pi}{6}+2\pi n\vee x=\frac{5\pi}{6}+2\pi n\vee x=\frac{7\pi}{6}+2\pi n \] where $n$ is an integer.

\section*{19}
\[ 2\sin^2 x-1=0 \]
\[ \sin ^2x=\frac{1}{2} \]
\[ \sin x=\pm\frac{1}{\sqrt{2}} \]
In the interval $0\le x\le 2\pi$, there are four solutions:
\[ x=\frac{\pi}{4}\vee x=2\pi-\frac{\pi}{4}=\frac{7\pi}{4}\vee x=\frac{\pi}{4}+\pi=\frac{5\pi}{4}\vee x=\pi-\frac{\pi}{4}=\frac{3\pi}{4} \]

\section*{20}
\[ 5\tan^2 x-15=0 \]
\[ \tan^2 x=3 \]
\[ \tan x=\pm\sqrt{3} \]
In the interval $0\le x\le 2\pi$, there are four solutions:
\[ x=\frac{\pi}{3}\vee x=2\pi-\frac{\pi}{3}=\frac{5\pi}{3}\vee x=\frac{\pi}{3}+\pi=\frac{4\pi}{3}\vee x=\pi-\frac{\pi}{3}=\frac{2\pi}{3} \]

\section*{21}
\[ 4\cos^2 x-1=0 \]
\[ \cos^2 x=\frac{1}{4} \]
\[ \cos x=\pm\frac{1}{2} \]
In the interval $0\le x\le 2\pi$, there are four solutions:
\[ x=\frac{\pi}{3}\vee x=2\pi-\frac{\pi}{3}=\frac{5\pi}{3}\vee x=\frac{\pi}{3}+\pi=\frac{4\pi}{3}\vee x=\pi-\frac{\pi}{3}=\frac{2\pi}{3} \] 

\section*{27}
\[ 2\tan^4x-\tan^2x-15=0 \]
This equation is in quadratic form with $u=\tan^2 x$. Thus we have
\[ \tan^2 x=\frac{1\pm \sqrt{(-1)^2-4(2)(-15)}}{2(2)} \]
\[ =\frac{1\pm \sqrt{121}}{4} \]
\[ =3\vee \pm -\frac{5}{2} \]
meaning \[ \tan x=\pm\sqrt{3}\vee \tan=\pm i\sqrt{\frac{5}{2}} \] The latter option is ruled out, as the range of $\tan x$ is only real numbers. Thus 
$\tan x=\pm\sqrt{3}$. In the interval $0\le x\le 2\pi$, there are four solutions:
As in problem 20, the solutions in the interval $0\le x\le 2\pi$ are:
\[ x=\frac{\pi}{3}\vee x=2\pi-\frac{\pi}{3}=\frac{5\pi}{3}\vee x=\frac{\pi}{3}+\pi=\frac{4\pi}{3}\vee x=\pi-\frac{\pi}{3}=\frac{2\pi}{3} \]
meaning the general solution of the equation is
\[ x=\frac{\pi}{3}+\pi n\vee x=\frac{5\pi}{3}+\pi n\vee x=\frac{4\pi}{3}+\pi n\vee x=\frac{2\pi}{3}+\pi n \] where $n$ is an integer.
which is the same as saying
\[ x=\frac{\pi}{3}+\pi n\vee x=\frac{2\pi}{3}+\pi n \]

\section*{28}
\[ \sqrt{\cos x}=2\cos x-1 \]
\[ \cos x=4\cos^2 x+1-4\cos x \]
\[ 4\cos^2 x-5\cos x+1=0 \]
This equation is in quadratic form with $u=\cos x$. Thus we have
\[ \cos x=\frac{5+\sqrt{(-5)^2-4(4)(1)}}{2(4)} \]
\[ =\frac{5+\sqrt{9}}{8} \]
\[ =1 \]
In the interval $0\le x\le 2\pi$, this holds when $x=0$ or $x=2\pi$. Thus the general solution to the equation is $x=2\pi n$ where $n$ is an integer.

\section*{29}
\[ 1+\cos x=\sqrt{3}\sin x \]
\[ (1+\cos x)^2=3\sin ^2x \]
\[ 1+\cos^2 x+2\cos x=3(1-\cos^2 x)\textrm{ Pythagorean identity} \]
\[ 1+\cos^2 x+2\cos x=3-3\cos^2 x \]
\[ 4\cos^2 x+2\cos x-2=0 \]
\[ (4\cos x-2)(\cos x+1)=0 \]
\[ \cos x=\frac{1}{2}\vee \cos x=-1 \]
In the interval $0\le x\le 2\pi$, this holds when 
\[ x=\frac{\pi}{3}\vee x=2\pi-\frac{\pi}{3}=\frac{5\pi}{3}\vee x=\pi \]
Now, we need to verify the above solutions:
\[ 1+\cos\frac{\pi}{3}\? \sqrt{3}\sin\frac{\pi}{3} \]
\[ 1+\frac{1}{2}\?\sqrt{3}\frac{\sqrt{3}}{2} \]
\[ \frac{3}{2}=\frac{3}{2} \] and
\[ 1+\cos\frac{5\pi}{3}\?\sqrt{3}\sin\frac{5\pi}{3} \]
\[ 1+\frac{1}{2}\?\sqrt{3}\times -\frac{\sqrt{3}}{2} \]
\[ \frac{3}{2}\ne -\frac{3}{2} \]
and
\[ 1+\cos\pi\?\sqrt{3}\sin\pi \]
\[ 1-1=\sqrt{3}\times 0 \]
meaning $\frac{5\pi}{3}$ is the only extraneous solution.
Thus the general solution to the equation is
\[ x=\frac{\pi}{3}+2\pi n\vee x=\pi +2\pi n \] where $n$ is an integer.

\section*{30}
\[ \sec x\csc^2 x=2\sec x \]
\[ \csc^2 x=2 \]
\[ \csc x=\pm\sqrt{2} \]
In the interval $0\le x\le 2\pi$, there are four solutions:
\[ \frac{\pi}{4},\pi+\frac{\pi}{4}=\frac{5\pi}{4},2\pi-\frac{\pi}{4}=\frac{7\pi}{4},\pi-\frac{\pi}{4}=\frac{3\pi}{4} \]
We now check the solutions:
\[ \sec\frac{\pi}{4}\csc^2\frac{\pi}{4}\? 2\sec\frac{\pi}{4} \]
\[ \sqrt{2}(\sqrt{2})^2\? 2\sqrt{2} \]
\[ 2\sqrt{2}=2\sqrt{2} \]
and
\[ \sec\frac{5\pi}{4}\csc^2\frac{5\pi}{4}\? 2\sec\frac{5\pi}{4} \]
\[ -\sqrt{2}(\sqrt{2})^2\? 2\times -\sqrt{2} \]
\[ -2\sqrt{2}=-2\sqrt{2} \]
and
\[ \sec\frac{7\pi}{4}\csc^2\frac{7\pi}{4}\? 2\sec\frac{7\pi}{4} \]
\[ \sqrt{2}(-\sqrt{2})^2\? 2\sqrt{2} \]
\[ 2\sqrt{2}=2\sqrt{2} \]
and
\[ \sec\frac{3\pi}{4}\csc^2\frac{3\pi}{4}\? 2\sec\frac{3\pi}{4} \]
\[ -\sqrt{2}(\sqrt{2})^2\? 2\times -\sqrt{2} \]
\[ -2\sqrt{2}=-2\sqrt{2} \]
meaning none of the apparent solutions are extraneous.

\section*{31}
\[ \sqrt{3}\cos^2 x=\cos^2 x\tan x \]
\[ \tan x=\sqrt{3} \]
In the interval $0\le x\le \pi$, the solution is $\frac{\pi}{3}$. We check it:
\[ \sqrt{3}\cos^2\frac{\pi}{3}\?\cos^2\frac{\pi}{3}\tan\frac{\pi}{3} \]
\[ \sqrt{3}\frac{1}{4}=\frac{1}{4}\sqrt{3} \]

\section*{32}
\[ 2\sin^2 x-\cos x-1=0 \]
\[ 2(1-\cos^2 x)-\cos x-1=0 \]
\[ -2\cos^2 x-\cos x+1=0 \]
\[ (-2\cos x+1)(\cos x+1)=0 \]
\[ \cos x=\frac{1}{2}\vee\cos x=-1 \]
The solutions in the interval $0\le x\le 2\pi$ are:
\[ \frac{\pi}{3},2\pi-\frac{\pi}{3}=\frac{5\pi}{3},\pi \]
We check them:
\[ 2\sin^2\frac{\pi}{3}-\cos\frac{\pi}{3}-1\? 0 \]
\[ 2\frac{3}{4}-\frac{1}{2}-1\? 0 \]
\[ \frac{3}{2}-\frac{3}{2}=0 \]
and
\[ 2\sin^2\frac{5\pi}{3}-\cos\frac{5\pi}{3}-1\? 0 \]
\[ 2\frac{3}{4}-\frac{1}{2}-1=0 \]
and
\[ 2\sin^2\pi-\cos\pi-1\? 0 \]
\[ 2(0)-(-1)-1\? 0 \]
\[ 0=0 \]
meaning none of the solutions are extraneous.

\section*{35}
\[ \cos x+\sin x\tan x=2 \]
\[ \cos x+\frac{\sin^2 x}{\cos x}=2\textrm{ Tangent Identity} \]
\[ \frac{\sin^2 x+\cos^2 x}{\cos x}=2 \]
\[ \frac{1}{\cos x}=2\textrm{ Pythagorean Identity} \]
\[ \cos x=\frac{1}{2} \]
In the interval $\pi\le x\le 2\pi$, the solution to this is $2\pi-\frac{\pi}{3}=\frac{5\pi}{3}$. We check it:
\[ \cos\frac{5\pi}{3}+\sin\frac{5\pi}{3}\tan\frac{5\pi}{3}\? 2 \]
\[ \frac{1}{2}+(-\frac{\sqrt{3}}{2})(-\sqrt{3})\? 2 \]
\[ \frac{1}{2}+\frac{3}{2}=2 \]
meaning the solution $x=\frac{5\pi}{3}$ is not extraneous.

\section*{36}
To find the points of intersection of the two given graphs in the interval $0\le x\le 2\pi$, we need to solve the equation
\[ 4\sin x+1=2\sin x+2 \]
\[ 2\sin x=1 \]
\[ \sin x=\frac{1}{2} \]
On the interval $0\le x\le 2\pi$, the solutions are $\frac{\pi}{6}$ and $\pi-\frac{\pi}{6}=\frac{5\pi}{6}$. We can check what these points are with this $x$-value:
\[ 4\sin\frac{\pi}{6}+1=3\wedge 4\sin\frac{5\pi}{6}+1=3 \] and 
\[ 2\sin\frac{\pi}{6}+2=3\wedge 2\sin\frac{5\pi}{6}+2=3 \] meaning $B$ is the correct answer.

\section*{39}
To find the points of intersection on the interval $0\le x\le 2\pi$, we nee d to solve the equation
\[ \sqrt{3}\tan^2 x=\sqrt{3}-2\tan x \]
\[ \sqrt{3}\tan^2 x+2\tan x-\sqrt{3}=0 \]
\[ (\sqrt{3}\tan x-1)(\tan x+\sqrt{3})=0 \]
\[ \tan x=\frac{1}{\sqrt{3}}\vee\tan x=-\sqrt{3} \]
In the interval $0\le x\le 2\pi$, the solutions to this are
\[ \frac{\pi}{6},\pi+\frac{\pi}{6}=\frac{7\pi}{6},\pi-\frac{\pi}{3}=\frac{2\pi}{3},2\pi-\frac{\pi}{3}=\frac{5\pi}{3} \]
We check the solutions (which will also tell us the $y$-coordinates of the points of intersection)
\[ \sqrt{3}\tan^2\frac{\pi}{6}\?\sqrt{3}-2\tan\frac{\pi}{6} \]
\[ \sqrt{3}\frac{1}{3}\?\sqrt{3}-2\frac{1}{\sqrt{3}} \]
\[ \frac{1}{\sqrt{3}}=\frac{1}{\sqrt{3}} \] (this also shows that $\frac{7\pi}{6}$ holds, as, because $\tan x$ has a period of $\pi$,
$\tan\frac{\pi}{6}=\tan\frac{7\pi}{6}$) and
\[ \sqrt{3}\tan^2\frac{2\pi}{3}\?\sqrt{3}-2\tan\frac{2\pi}{3} \]
\[ 3\sqrt{3}\?\sqrt{3}-2(-\sqrt{3}) \]
\[ 3\sqrt{3}=3\sqrt{3} \] (this also shows that $\frac{5\pi}{3}$ holds) meaning the points of intersection on the interval $0\le x\le 2\pi$ are
\[ (\frac{\pi}{6},\frac{1}{\sqrt{3}}),(\frac{7\pi}{6},\frac{1}{\sqrt{3}}),(\frac{2\pi}{3},3\sqrt{3}),(\frac{5\pi}{3},3\sqrt{3}) \]

\section*{45}
(a): We are given the point $(2,3)$, meaning $x=2$ and $y=3$, and that $\theta=\frac{\pi}{3}$. We know that 
\[ x'=2\cos\frac{\pi}{3}-3\sin\frac{\pi}{3} \]
\[ =2\frac{1}{2}-3\frac{\sqrt{3}}{2} \]
\[ =\frac{2-3\sqrt{3}}{2} \]
\[ y'=2\sin\frac{\pi}{3}+\cos\frac{\pi}{3} \]
\[ =2\frac{\sqrt{3}}{2}+\frac{1}{2} \]
\[ =\frac{2\sqrt{3}+2}{2} \]


(b): We are given the point $(6,2)$, and that it has been rotated to $(3\sqrt{3}-1,\sqrt{3}+3)$. Thus we know that
\[ 3\sqrt{3}-1=6\cos\theta-2\sin\theta \]
and
\[ \sqrt{3}+3=6\sin\theta+2\cos\theta \]
\[ 2\cos\theta=\sqrt{3}+3-6\sin\theta \]
and thus
\[ 3\sqrt{3}-1=3(2\cos\theta)-2\sin\theta \]
\[ 3\sqrt{3}-1=3(\sqrt{3}+3-6\sin\theta)-2\sin\theta \]
\[ 3\sqrt{3}-1=3\sqrt{3}+9-18\sin\theta-2\sin\theta \]
\[ -10=-20\sin\theta \]
\[ \sin\theta=\frac{1}{2} \] and thus $\theta=\frac{\pi}{6}$.
We check to make sure this angle realy works:
\[ \sqrt{3}+3\? 6\sin\frac{\pi}{6}+2\cos\frac{\pi}{6} \]
\[ \sqrt{3}+3\? 3+2\frac{\sqrt{3}}{2} \]
\[ \sqrt{3}+3=\sqrt{3}+3 \]
and
\[ 3\sqrt{3}-1\? 6\cos\frac{\pi}{6}-2\sin\frac{\pi}{2} \]
\[ 3\sqrt{3}-1\? 6\frac{\sqrt{3}}{2}-1 \]
\[ 3\sqrt{3}-1=3\sqrt{3}-1 \]
meaning $\theta=\frac{\pi}{6}$ is the correct solution.

\end{document}

